\section{Uniform contact continuation across physical clearance seams}
\label{sec:v66-contact-continuation}
The protected collars of Section~\ref{sec:v45-critical-collars} keep
both the selected contact equations and every physical decision fixed.
For a clearance caustic the latter requirement is inappropriate: the
physical source occupies only part of a smooth selected-contact chart.
We separate the two requirements. The chart below continues selected
reflections, whereas all first-hit and section restrictions will remain
as indicators in the original source integral.

Set $q=47/53$, $\sigma=\sqrt q$, $d_j=\min(j,m-j)$, and
$e_j=q^{d_j/4}$. A selected normal-to-normal word is
\emph{$\chi$-incident} if it is in the physical-side closure, its
selected contacts have positive incidence, and
\begin{equation}\label{eq:v66-tapered-incidence}
 c_j(z)\ge\chi e_j\quad(0\le j\le m),\qquad 0<\chi\le1/10.
\end{equation}
There is no clearance or section-margin condition in this definition.
In particular several competing disks and section edges may meet at $z$.
All constants in this section depend only on the fixed one-flight
geometry and contact charts, not on $m,R,z$, or their exact labels.

\subsection{An incidence-normalized inverse}
Use the oriented contact arclengths and the matrix $A$ of
\eqref{eq:v45-contact-matrix}. Write $D_c=\operatorname{diag}(c_j)$.
The following refinement of the inherited Neumann estimate is useful.

\begin{lemma}[Weighted contact inverse]\label{lem:v66-weighted-inverse}
For every positive selected incidence,
\begin{equation}\label{eq:v66-refined-inverse}
 |(A^{-1})_{ij}|\le C\min(c_i,c_j)q^{|i-j|},\qquad
 |(D_c^{-1}A^{-1})_{ij}|\le Cq^{|i-j|}.
\end{equation}
Consequently $N=A D_c$, the internal Hessian with row $j$ divided
by $c_j$, has an inverse bounded independently of $m$ in the norm
\[
 \|v\|_\sigma=\max_{1\le j<m}\sigma^{-d_j}|v_j|.
\]
The assertions about internal matrices are vacuous for $m=1$.
\end{lemma}
\begin{proof}
Write $A=D(I-K)$, where $D$ is its diagonal. The row sums of $K$
are at most $q$, its entries are nonnegative, and $K$ connects only
neighboring indices. Thus $(K^r)_{ij}=0$ for $r<|i-j|$ and
$|(K^r)_{ij}|\le q^r$. Moreover $D_{jj}^{-1}\le Rc_j/2$.
The Neumann series gives
$|(A^{-1})_{ij}|\le (R_+/2)(1-q)^{-1}c_jq^{|i-j|}$.
Symmetry gives the same inequality with $c_i$. Division by $c_i$
proves the second assertion. Since $|d_i-d_j|\le|i-j|$,
\[
 \sup_i\sum_j q^{|i-j|}\sigma^{d_j-d_i}
 \le 1+2\sum_{r\ge1}(q/\sigma)^r<\infty.
\]
This is the asserted weighted inverse bound. No product of inverse
incidences remains in it.
\end{proof}

\begin{theorem}[Count-uniform selected-contact jets]
\label{thm:v66-contact-jets}
There are $a_0>0$ and $C_1,C_2,C_3<\infty$ such that every
$\chi$-incident word has a unique selected-contact continuation on
$|u|,|v|<a_0\chi^3$ about its normal endpoints. On this square all
selected incidences are at least $\chi e_j/2$. If $y_j(u,v)$ is
internal arclength displacement from the original contact, then
\begin{equation}\label{eq:v66-contact-jets}
 \|D^r y_j\|\le C_r\chi^{1-r}\sigma^{d_j},
              \qquad 1\le r\le3.
\end{equation}
The same bounds, with fixed changes of constants and neighboring
indices, apply to the selected positions, velocities and flight data.
The continuation need not satisfy the first-hit or section tests.
Only its restriction satisfying those tests represents physical source.
\end{theorem}
\begin{proof}
Let $E_j$ be the internal stationary equation for the polygonal length.
Freeze $c_j^0=c_j(z)$ and put $G_j=E_j/c_j^0$. At the base point
$D_yG=D_c^{-1}(D_cAD_c)=A D_c=N$. Its inverse has the weighted
bound of Lemma~\ref{lem:v66-weighted-inverse}. The endpoint forcing
has only its first and last rows nonzero. Their entries have modulus
$c_0t_0$ and $c_mt_{m-1}$ after the indicated row division, and
therefore have bounded weighted norm.

Each $E_j$ depends on three adjacent contact positions. Its derivatives
of any fixed order are bounded on fixed local arclength charts: the
formulas divide by neighboring flight lengths, which are at least
$\ell_0=3/50$, and not by incidence. Tangent signs are fixed from the
base contact word. Put $\omega_j=\sigma^{d_j}$, with endpoint weights
one. Neighboring weights differ by at most $\sigma^{-1}$. For $r\ge2$
a derivative of order $r$, with arguments of weighted norm at most one,
therefore satisfies
\begin{equation}\label{eq:v66-nonlinear-weight}
 \omega_j^{-1}|D^rG_j[v_1,\ldots,v_r]|
 \le C_r\frac{\omega_j^{r-1}}{c_j^0}
 \le C_r\chi^{-1}q^{((r-1)/2-1/4)d_j}
 \le C_r\chi^{-1}.
\end{equation}
The same estimate holds for mixed endpoint derivatives, whose support
is confined to the endpoint rows. These are estimates for the scaled
norms in varying finite dimensions; their constants do not count rows.

Apply the contraction map obtained by multiplying the stationary
equations by the fixed inverse $N^{-1}$. On a weighted ball of radius
$a\chi$ and endpoint square of radius $a\chi$, the derivative of its
nonlinear part is bounded by $C a$. Taking $a$ small gives a contraction
into a slightly larger fixed multiple of the endpoint radius. This
proves a unique smooth solution and a uniformly bounded inverse
linearization there. In particular its first derivative has weighted
norm at most $C$. Differentiating the equation twice gives a second
derivative bounded by $C\chi^{-1}$; differentiating a third time
gives $C(\chi^{-1}+\chi^{-2})\le C'\chi^{-2}$, by
\eqref{eq:v66-nonlinear-weight}. This proves
\eqref{eq:v66-contact-jets}, on the smaller displayed square.

Positions and unit flight directions are Lipschitz functions of
adjacent arclengths with uniform constants. Thus incidence changes by
at most $C(|u|+|v|)\omega_j$. Relative to $c_j^0$ this is at most
$C\chi^{-1}(|u|+|v|)q^{d_j/4}$. Reduction of $a_0$ keeps it below
$1/2$. Endpoint incidences likewise stay above $1/2$. The local
stationary solution thus retains the selected positive reflection
signs. The constants were chosen inside fixed contact charts.
Uniqueness on overlaps gives the whole square. This argument used
neither a lower bound on a competing-hit margin nor a fixed section
membership. On points satisfying the original first-hit tests it
agrees with the original billiard word by endpoint injectivity in
Lemma~\ref{lem:v44-endpoint-curvature}.
\end{proof}

\subsection{Uniform Morse coordinates and relative density}
Let $F_z$ be the selected reduced action and let
$w_z=|\partial_u\partial_vF_z|/(4\pi R c_*)$, where
$c_*=91/(10000\pi)$. This is a smooth coefficient on the continuation;
it is the density of the original source only after imposing its
physical and section indicators.

\begin{theorem}[A uniform quadratic chart at the seam]
\label{thm:v66-uniform-morse}
There are $r_0,C>0$ such that, putting
\begin{equation}\label{eq:v66-morse-budget}
 r_\chi=r_0\chi^3,\qquad h_\chi=r_\chi^2/2,
                 \qquad K_\chi=C\chi^{-2},
\end{equation}
every $\chi$-incident word has a selected-contact coordinate map
$x=\Phi_z(y)$ on $|y|<r_\chi$ satisfying
\begin{align}\label{eq:v66-uniform-morse}
 F_z(\Phi_z(y))&=t_z+|y|^2/2, &
 |p_0|+|p_m|&\le C|y|,\notag\\
 \widetilde w_z(y)&=w_z(\Phi_z(y))|\det D\Phi_z(y)|,&
 \widetilde w_z(0)&=J_z/(2\pi),\notag\\
 \left|\log\frac{2\pi\widetilde w_z(y)}{J_z}\right|
                    &\le K_\chi|y|.
\end{align}
Here $J_z$ is the unchanged reversible coefficient of
Theorem~\ref{thm:v65-reversible-weight}, interpreted from its physical
side. The constants are uniform in the word length and in every
positive incidence satisfying \eqref{eq:v66-tapered-incidence}.
The disk is a coordinate disk, not an assertion of a full physical disk.
\end{theorem}
\begin{proof}
The selected contact second variation has the same positive internal
block throughout the continuation. Its endpoint Schur complement
therefore gives
\[
 \lambda I\le \nabla^2F_z\le\Lambda I,
 \quad \lambda=1/(2R_+),\quad
 \Lambda=2(1/R_-+1/\ell_0),
\]
after reducing the endpoint square. First variation expresses $DF_z$
using only the two endpoint flight directions. The first three
contact jets consequently bound $D^rF_z$, $2\le r\le4$, by
$C\chi^{-2}$. The value $F_z(0)$, which grows with $m$, is not
included in this derivative estimate.

The determinant formula \eqref{eq:v61-cross-and-product} is valid on
the selected continuation. In its logarithmic derivative the flight
terms have summable bounds $C\omega_j$. The diagonal potential terms
of $\operatorname{tr}(A^{-1}dA)$ are bounded, using the refined diagonal
inverse, by $C|dc_j|/c_j$. Their sum is at most
\[
 C\chi^{-1}\sum_{j=1}^{m-1}q^{d_j/4}|(du,dv)|
                    \le C'\chi^{-1}|(du,dv)|.
\]
Off-diagonal terms and the two endpoint cosines have bounded sums.
It follows that $|D\log w_z|\le C\chi^{-1}$. This is a relative
bound and does not require a lower bound for $w_z(0)$ as $m$ grows.

For completeness use the explicit quadratic construction
\[
 Q_z(x)=2\int_0^1(1-s)\nabla^2F_z(sx)\,ds,
                  \qquad y=Q_z(x)^{1/2}x.
\]
It has derivative $(\nabla^2F_z(0))^{1/2}$ at zero. Matrix square
root and inversion have bounded derivatives on matrices with spectrum
in $[\lambda,\Lambda]$. The preceding fourth-derivative bound gives
$\|D^2(y(x))\|\le C\chi^{-2}$ on the endpoint square. On a square
of radius $a_0\chi^3$ its first derivative is uniformly close to its
invertible value at zero, after reducing $a_0$. The quantitative
inverse function argument supplies a ball $|y|<r_0\chi^3$ in its
image, with $D\Phi_z$ bounded and $D^2\Phi_z$ bounded by
$C\chi^{-2}$. The coordinate identity is exact. Differentiating
$\log\widetilde w_z=\log(w_z\circ\Phi_z)+\log|\det D\Phi_z|$
gives a bound $C\chi^{-2}$. Integration from zero proves the last
line of \eqref{eq:v66-uniform-morse}. At zero its value is the original
Morse coefficient $J_z/(2\pi)$. The momentum estimate follows from
first variation, the upper Hessian bound and $|\Phi_z(y)|\le C|y|$.
\end{proof}

\paragraph{Coverage provided by the theorem.}
Every normal critical point in the physical-side closure satisfying
\eqref{eq:v66-tapered-incidence} has this same chart budget. No
nonzero clearance margin, section margin, individual boundary-gradient
lower bound, or angle between boundary gradients enters the budget.
Selected grazing and points violating the tapered incidence envelope
are not included. The next section keeps their source, and every
noncritical or uncovered source component, in an explicit positive
complement. No physical probability is assigned to the continued graph.
